\documentclass[10pt,a4paper]{amsart}
\usepackage[margin=19mm]{geometry}
\usepackage{amsmath,amssymb,mathtools}
\usepackage[scr=rsfs,cal=euler]{mathalfa}
\usepackage{xfrac}
\usepackage[unicode,hidelinks,bookmarks=false]{hyperref}
\newcommand{\ie}{i.e.\ }
\newcommand{\cf}{cf.\ }
\newcommand{\resp}{respectively\ }
\newcommand{\Subsection}{\subsection}
\providecommand{\nobreakdash}{\nobreak\hbox{-}}
\setlength{\emergencystretch}{2em}
\multlinegap0pt
\usepackage{bm}
\usepackage{bbm}
\usepackage{dsfont}
\usepackage{xspace}
\usepackage{cancel}
\usepackage{mathtools}
\usepackage{amsthm}
\makeatletter
\@ifundefined{theorem}{\newtheorem{theorem}{Theorem}}{}
\makeatother
\title[Borel 1 type mappings and the respective equi-families]{Borel 1 type mappings and the respective equi-families}
\author{Marek Balcerzak, \v{L}ubica Hol\'a, Olena Karlova, Piotr Szuca}
\date{Source: arXiv:2512.15912v3 (CC BY 4.0)}
\begin{document}
\maketitle

\noindent\emph{Source and licence.} Borel 1 type mappings and the respective equi-families. Version arXiv:2512.15912v3 (CC BY 4.0), distributed under the Creative Commons Attribution 4.0 International licence, as recorded in the arXiv licence metadata for this exact version and shown on the abstract page https://arxiv.org/abs/2512.15912v3. Full rights notice: licenses/arxiv-2512.15912-CC-BY-4.0.txt.\par\smallskip

\noindent\textit{Editorial scope.} Counted unit: the Theorem whose source label is \texttt{None} in the article (source0.tex lines 896--905, source label \texttt{None}), delivered together with the article's own complete original proof of it (source0.tex lines 907--946, one \texttt{proof} environment of 40 lines). No mathematics is written, repaired or re-derived: statement and proof are the article's own text line for line. Required context retained uncounted: none. Every definition, display and label of those lines is retained unchanged. Omitted from this package and not redistributed: 1--895, 906--906, 947--975. Nothing inside a retained excerpt is omitted or altered: every retained body is delivered exactly as the article's lines are written, so no omission receipt of any kind accompanies this package. Citations: the retained excerpts carry no citation command at all, so no bibliography entry of the article is transcribed and no reference is redistributed. Reference adaptation: none was needed and none was made; every reference inside the delivered unit resolves inside it, so no unresolved reference or citation is produced. Printed slips kept literal: no printed word, symbol, hypothesis or number of the counted statement or of its proof was altered. Layout and class: the article's own class is \texttt{amsart} with packages including extsizes, amsfonts, amsmath, amsthm, amscd, amssymb, latexsym, mathrsfs, enumitem, tikz, babel, color, geometry; the wrapper is the lane's pinned \texttt{amsart} editorial wrapper at 10pt on A4, which restates the theorem-like environments the retained text needs and adds the article's own macro definitions that the retained text needs (none), with extra wrapper packages (bm, bbm, dsfont, xspace, cancel, mathtools). The delivered numbering is the unit's own. Assets: no retained span contains a figure environment, a graphics-inclusion command, a PDF-page inclusion, a table or a TikZ picture; no image file is copied into this package, the package declares no asset dependency and no image file is redistributed with it. Licence: version arXiv:2512.15912v3 of this article is distributed under the Creative Commons Attribution 4.0 International licence, as recorded in the arXiv licence metadata for this exact version and shown on the abstract page https://arxiv.org/abs/2512.15912v3. A full rights notice is delivered at licenses/arxiv-2512.15912-CC-BY-4.0.txt. Characters: every retained excerpt is pure ASCII, so the delivered engine, XeLaTeX with its default Latin Modern text font, reports no missing character.
 \begin{theorem}
 Let $X, Y$ and $(Z,d)$ be metric spaces, $1\le\alpha<\omega_1$ and let a function $f\colon X\times Y\to Z$ satisfy the following properties:
 \begin{enumerate}
 	\item the family $\mathscr F_X=\{f^x: x\in X\}$ is equi-$\alpha$-GLP,
 	
 	\item for every $y\in Y$ the section $f_y\colon X\to Z$ belongs to the Borel class $\alpha$. 
 \end{enumerate}
 
 Then $f\colon X\times Y\to Z$ belongs to the  Borel class $\alpha$.
\end{theorem}

\begin{proof} We will show that for every $\varepsilon>0$ there exists a function $g\colon X\times Y\to Z$ of Borel class $\alpha$ such that
	\begin{center}
		 $d(f(p),g(p))\le\varepsilon$ for all $p\in X\times Y$.
	\end{center} 
	Fix $\varepsilon>0$. Since $\mathscr F_X$ is equi-$\alpha$-GLP, there exists a $\sigma$-discrete family $\mathcal A$ in $Y$ consisting of sets of additive and multiplicative class $\alpha$ simultaneously such that
	\begin{center}
		  $Y=\bigcup \mathcal A$ \quad and \quad ${\rm diam}\,f^x(A)\le\varepsilon$
	\end{center} 
	for all $A\in\mathcal A$ and $x\in X$. Let $\mathcal A=\bigcup\limits_n\mathcal A_n$, where each family $\mathcal A_n$ is discrete in $Y$. For every $n\in\mathbb N$ we put 
	\begin{center}
		 	$\mathcal B_n=\{ A\setminus \bigcup\limits_{k<n}\bigcup \mathcal A_k: A\in\mathcal A_n\}$\,\,\,and\,\,\,$B_n=\bigcup\{B:B\in\mathcal B_n\}$.
	\end{center}
    Notice that each family $\mathcal B_n$ is discrete (in particular, disjoint), since $\mathcal B_n\prec\mathcal A_n$.  Let 
      $$
      \mathcal B=\bigcup\limits_{n=1}^\infty \mathcal B_n.
      $$
 	Then $\mathcal B$ is a  $\sigma$-discrete family of sets of ambiguous class $\alpha$ with $\mathcal B\prec \mathcal A$ and $\bigcup\mathcal B=Y$. Moreover, $B'\cap B''=\emptyset$ for all $B'\in\mathcal B_k$, $B''\in\mathcal B_n$ and distinct $k,n\in\mathbb N$. Notice that every set $B_n$ belongs to the additive and multiplicative class $\alpha$ simultaneously as the union of a discrete family of sets of the same classes in metric space $Y$. 
	
	Take an arbitrary $y_{B,n}\in B$ for every nonempty $B\in\mathcal B_n$, $n\in \mathbb N$. For each $y\in Y$ we put $n(y)=\min\{n\in\mathbb N: y\in B_n\}$. Now let $(x,y)\in X\times Y$. Since the family $\mathcal B_{n(y)}$ is disjoint, there exists a unique set $B(y)\in\mathcal B_{n(y)}$ such that $y\in B(y)$. We define
	$$
	g(x,y)=f(x,y_{B(y),n(y)}).
	$$
	Then for every point $(x,y)\in X\times Y$ we have 
	$$
	d(f(x,y),g(x,y))=d(f(x,y), f(x,y_{B(y),n(y)}))\le {\rm diam}\, f^x(B(y)),
	$$
	since $y, y_{B(y),n(y)}\in B(y)$. It follows from the property $\mathcal B\prec\mathcal A$ that ${\rm diam}\, f^x(B(y))\le {\rm diam}\, f^x(A)$ for some $A\in\mathcal A$. Hence,
	$$
		d(f(x,y),g(x,y))\le\varepsilon.
	$$

It remains to prove that $g\colon X\times Y\to Z$ belongs to the Borel class $\alpha$. Fix $n\in\mathbb N$ and let 
$g_n=g\restriction_{X\times B_n}$. It follows from the definition of $g$ and assumption~(2) of the theorem, that $g_{n,B}=g_n\restriction_{X\times B}$ belongs to the Borel class $\alpha$ for every $B\in\mathcal B_n$. Then for every open set $W\subseteq Z$ we have that the preimage
$$
g_n^{-1}(W)=\bigcup\{g_{n,B}^{-1}(W):B\in\mathcal B_n\}
$$
is of the additive class $\alpha$ in $X\times B_n$ as a union of a discrete family of sets of the same class. Now, since $\{X\times B_n:n\in\mathbb N\}$ is a countable covering of $X\times Y$ by sets of ambiguous class $\alpha$ and $g\restriction_{X\times B_n}$ is Borel $\alpha$, we conclude that $g\colon X\times Y\to Z$ belongs to the Borel class $\alpha$.

	Finally, let us remember that the family of functions of Borel class $\alpha$ between metric spaces is closed under uniform limits.	 
\end{proof}

\end{document}
